%==================================================
%----- 2.9 Condições iniciais e cenários de validação
%==================================================
\section{Condições iniciais e cenários de validação}

Esta seção fixa o estado inicial, as referências de atitude e os ensaios empregados nas simulações.
As grandezas respeitam as restrições holonômicas e suas derivadas no instante inicial, conforme
\eqref{eq:ci_sist_linear_blocos1} a
\eqref{eq:ci_sist_linear_blocos2},
e a dinâmica dada em
\eqref{eq:dinamica_mult_lagrange}, 
ou sua forma com dissipação em \eqref{eq:dae_diss}.

\subsection{Condições iniciais}
O estado inicial é
$\vec x(0)=\begin{bmatrix}\vec q(0) & \dot{\vec q}(0)\end{bmatrix}^{\!\top}$, com:
\begin{itemize}
\item Consistência dos vínculos: $\Phi(\vec q(0),t_0)=\vec 0$ e $\mathbf J(\vec q(0),t_0)\,\dot{\vec q}(0)+\tfrac{\partial\Phi}{\partial t}\big|_{(\vec q(0),t_0)}=\vec 0$.
\item Atitude da base: orientação $R_{B}^{I}(0)$ (ou quaternion $q_b(0)$, normalizado) e velocidade angular $\boldsymbol\omega_b(0)$.
\item Manipulador robótico: juntas $\vec q_m(0)$ e $\dot{\vec q}_m(0)$ dentro dos limites.
\item Microgravidade: $V(\vec q)\approx 0$ e $\mathbf g(\vec q)=\vec 0$.
\item Dissipação: quando presente, $\mathbf D\succeq \mathbf 0$ conforme
\eqref{eq:Rayleigh} e \eqref{eq:Qd}.
\end{itemize}

\subsection{Referências de atitude}
As referências de atitude e velocidade são definidas conforme o cenário:
\begin{enumerate}[label=(\roman*)]
\item referência constante (apontamento fixo);
\item referência nula (controlador de atitude desligado);
\item trajetória suave de atitude para casos com controle ativado.
\end{enumerate}

\subsection{Cenários de validação}
Dois ensaios compõem a validação:
o primeiro, é a atitude sem controle, onde o manipulador robótico permanecerá travado na posição inicial;
o controle de atitude está desligado,
sendo $\vec q_m(t)\equiv \vec q_m(0)$.
Os gráficos obtidos evidenciam o comportamento da dinâmica livre da base.
Enquanto que o segundo cenário mantém-se a atitude sem controle, mas desta vez, o manipulador robótico está realizando movimento, além disso, o controle de atitude permanece desligado, com perfil $\vec q_m(t)$.
Neste cenário, as reações base--manipulador tornam-se evidentes.


\subsection{Sinais e métricas}
As grandezas acompanhadas incluem:
\begin{itemize}
\item atitude da base (quaternion $q_b(t)$ ou parametrização mínima) e $\boldsymbol\omega_b(t)$; erro de atitude quando houver referência;
\item torques/forças generalizadas $\vec\tau(t)$, quando pertinentes;
\item momento angular da base $H_b(t)$ e variações associadas;
\item consistência dos vínculos: $\|\Phi(\vec q(t),t)\|$ e $\big\|\mathbf J(\vec q,t)\dot{\vec q}+\tfrac{\partial \Phi}{\partial t}\big\|$;
\item energia mecânica $E(t)$ e taxa $\dot E(t)=-\dot{\vec q}^{\top}\mathbf D\,\dot{\vec q}$ quando
\eqref{eq:Rayleigh} é empregado.
\end{itemize}

\subsection{Reprodutibilidade}
Para cada ensaio ficam explicitados:
$(\vec q(0),\dot{\vec q}(0))$ completo (base e juntas),
os parâmetros inerciais e a convenção de quaternion,
a presença/ausência de $\mathbf D$,
bem como os ganhos de estabilização de Baumgarte
$(\zeta,\omega_n)$ quando aplicável,
conforme \eqref{eq:acclevel_baumgarte}.
